% !TEX root = cha-approach.tex
\documentclass[../../main.tex]{subfiles}

\begin{document}

% \section{Summary}

% This chapter presented the approach implemented in \wassis to generate software visualizations by combining Multidimensional Scaling and Wave Function Collapse. The approach consists of four main phases (\figref{fig:approach_flow}). First, the source code is parsed to extract classes and their relationships, and an empty grid is initialized. Second, MDS is used to place classes on the grid according to the distances between them. The distances are computed from the classes' relationships, and the Floyd-Warshall algorithm is used to consider indirect relationships. Classes can occupy multiple cells when they have many connections, providing enough space for their paths. Third, the relationships between classes are represented as paths and generated using WFC. The generated paths can be optionally improved using Dijkstra's algorithm. Finally, quantitative metrics are computed and the resulting visualization is rendered.

% To reduce visual clutter, relationships with the same destination and relationship type are merged into a single path. Each merged path is assigned a unique path ID, which is propagated through the grid to ensure that paths connect the intended classes. The observation strategy selects cells with the highest entropy, allowing path cells to be generated before empty cells. Tile weights are also used to avoid unnecessary crossings and junctions.

% Several extensions were introduced to make the approach tractable for larger software systems. A placeholder tile system prevents the number of tiles from growing with the number of paths, while a dedicated path ID propagation step resolves path identities during generation. Directional path exits and dynamic tile weighting guide paths towards their targets. In addition, backtracking strategies allow WFC to recover from contradictions by changing previous decisions, restarting the generation, or expanding the grid when necessary. 

\section{Summary}

This chapter presented an approach for visualizing software systems by combining Multidimensional Scaling and Wave Function Collapse. MDS provides an initial layout by positioning classes according to their relationships, while WFC generates paths between them on a grid. The main goal is to preserve the connectivity of the software system while producing clear visualizations with limited overlaps and unnecessary crossings.

Adapting WFC to this problem requires more than defining tiles and their local constraints. The approach must also ensure that paths connect the intended classes and that the generation process prioritizes the connections that need to be drawn. To achieve this, relationships are merged when they share the same destination and type, and the observation strategy favors cells that contribute to path generation. Additional mechanisms guide paths towards their targets and help the algorithm recover when its decisions lead to contradictions.

These ideas define the overall approach, but their practical application also depends on how the different components interact and how the algorithm manages the generation process. The next chapter describes the implementation of the approach, including its main components and the mechanisms used to put it into practice.

\end{document}